%! The Tangent Function — Unit 3, Lesson 3.8 — Exit Ticket
\documentclass[10pt]{article}
\usepackage{precalculus-article}
\usepackage{precalculus-boxes}
\newcommand{\TallMath}[1]{$\displaystyle #1\rule[-1.4em]{0pt}{3.2em}$}

\begin{document}

\pageheader{Unit 3, Lesson 3.8}{Exit Ticket}
\namedateperiod

\vspace{0.10in}

\begin{notesbox}{Exit Ticket --- Answer independently}

\textbf{1.}\quad Given $\sin\theta = \dfrac{\sqrt{3}}{2}$ and $\cos\theta = \dfrac{1}{2}$,
compute $\tan\theta$.  Show your work.

\vspace{4pt}
$\tan\theta = \dfrac{\sin\theta}{\cos\theta} = $ \writeline

\vspace{0.14in}
\textbf{2.}\quad State the period and the location of the two positive vertical
asymptotes closest to $\theta = 0$ for the function $g(\theta) = \tan(2\theta)$.

\vspace{4pt}
Period of $g$: \writeline

\vspace{4pt}
Asymptotes (two positive values closest to 0): $\theta = $ \blank{2.2cm} and $\theta = $ \blank{2.2cm}

\vspace{0.14in}
\textbf{3.}\quad For the parent function $y = \tan\theta$, the graph is always
\underline{\hspace{2.2cm}} (increasing / decreasing) between consecutive
asymptotes, and changes from concave \underline{\hspace{2cm}} to concave
\underline{\hspace{2cm}} within each period.

\end{notesbox}

\end{document}
